%%%%%%%%%%%%%%%%%%%%%%%%%%%%%%%%%%%%%%%%%
% Główny plik pracy
% Szablon pracy dyplomowej
% Wydział Informatyki 
% Zachodniopomorski Uniwersytet Technologiczny w Szczecinie
% autor Joanna Kołodziejczyk (jkolodziejczyk@zut.edu.pl)
% Bardzo wczesnym pierwowzorem szablonu był
% The Legrand Orange Book
% Version 2.1 (26/09/2018)
%
% Modifications to LOB assigned by %JK
%%%%%%%%%%%%%%%%%%%%%%%%%%%%%%%%%%%%%%%%%


%----------------------------------------------------------------------------------------
%	PAKIETY ORAZ PLIKI ZAWIERAJĄCE DEFINICJE STYLI I KONFIGURACJĘ OTOCZEŃ LATEX
%----------------------------------------------------------------------------------------

\documentclass[12pt,fleqn,twoside]{book} % JK Rozmiary czcionek są zgodne z wymaganiami uczelni 
% UWAGA! - wydruk jest dwustronny i jest to zabieg zamierzony

%%%%%%%%%%%%%%%%%%%%%%%%%%%%%%%%%%%%%%%%%
% Plik konfigurujący
% Szablon pracy dyplomowej
% Wydział Informatyki 
% Zachodniopomorski Uniwersytet Technologiczny w Szczecinie
% autor Joanna Kołodziejczyk (jkolodziejczyk@zut.edu.pl)
% Bardzo wczesnym pierwowzorem szablonu był
% The Legrand Orange Book
% Version 2.1 (26/09/2018)
%
% Modifications to LOB assigned by %JK
%%%%%%%%%%%%%%%%%%%%%%%%%%%%%%%%%%%%%%%%%


%----------------------------------------------------------------------------------------
%	VARIOUS REQUIRED PACKAGES AND CONFIGURATIONS
%----------------------------------------------------------------------------------------

%G EOMETRY
\usepackage[top=3.2cm,bottom=3.5cm,left=2.5cm,right=2.5cm,headsep=1.5ex,bindingoffset=1cm,a4paper]{geometry} % Page margins

% GRAPHICS
\usepackage{graphicx} % Required for including pictures
\graphicspath{{Pictures/}} % Specifies the directory where pictures are stored

% For math equations, theorems, symbols, etc
\usepackage{amsmath,amsfonts,amssymb,amsthm} 

% Customize lists
\usepackage{enumitem} 
\setlist{nolistsep} % Reduce spacing between bullet points and numbered lists
\setlist[itemize]{label=--}

% Required for nicer horizontal rules in tables
\usepackage{booktabs} 

\usepackage{xcolor} % Required for specifying colors by name
% Define blue colors used for highlighting throughout the book- based on the WI ZUT colors
\definecolor{blueWI}{cmyk}{.6,.2,0,.0} % JK - Define the blue colour used for highlighting throughout the book
\definecolor{blueZUT}{cmyk}{1,.75,0,.2} % JK - Define the blue colour used for highlighting throughout the book
\definecolor{grayZUT}{cmyk}{0,0,0,0.4} % JK - Define the blue colour used for highlighting throughout the book

%other definision
\newcommand{\rulecolor}[1]{\color{#1}\rule}

%----------------------------------------------------------------------------------------
%	FONTS AND LANGUAGE (JK - configuring and styling)
%----------------------------------------------------------------------------------------
\usepackage[sfdefault, book,]{FiraSans} %% option 'sfdefault' activates Fira Sans as the default text font
\usepackage[T1]{fontenc}
\renewcommand*\oldstylenums[1]{{\firaoldstyle #1}}

\usepackage{newtxmath,newtxtext}
% \usepackage{t1enc}
\usepackage[polish]{babel}% moved after to avoid conflict between polish babel and amsmath
\usepackage[utf8]{inputenc} 
%\usepackage{avant} % Use the Avantgarde font for headings

%----------------------------------------------------------------------------------------
%	PAGE HEADERS AND FOOTERS (JK - Styling for the current chapter in the header)
%----------------------------------------------------------------------------------------

\usepackage{fancyhdr} % Required for header and footer configuration
\setlength{\headheight}{2.5ex}
\pagestyle{fancy} % Enable the custom headers and footers

\renewcommand{\chaptermark}[1]{\markboth{\sffamily\normalsize\thechapter.\hspace{5pt} #1}{}} % JK - Styling for the current chapter in the header
\renewcommand{\sectionmark}[1]{\markright{\sffamily\normalsize\thesection\hspace{5pt} #1}{}} % Styling for the current section in the header

\fancyhf{} % Clear default headers and footers

% JK - header with page hanging  and chapter title in the box
\fancyhead[EL]{%
  \textcolor{white}{%
    \llap{%
      \colorbox{blueZUT}{%
        \makebox[7ex][r]{\sffamily\thepage}%
      }%
      \hspace{1.25\marginparsep}%
      \hspace{-\fboxsep}%
    }%
    \textcolor{black}{%
      \colorbox{blueWI!20}{%
        \makebox[\textwidth][l]{\ttitle}}}
     \addtolength{\headheight}{5ex} % Increase the spacing around the header slightly
  }%
}
\fancyhead[OR]{%
  \textcolor{white}{%
    \llap{%
      \colorbox{blueZUT}{%
        \makebox[7ex][r]{\sffamily\thepage}%
      }%
      \hspace{1.25\marginparsep}%
      \hspace{-\fboxsep}%
    }%
    \textcolor{black}{%
      \colorbox{blueWI!20}{%
        \makebox[\textwidth][l]{\leftmark}}}%{\rightmark}}}
    \addtolength{\headheight}{5ex} % Increase the spacing around the header slightly    
  }%
}

\fancypagestyle{plain}{%
   \fancyhead{} % get rid of headers
   \fancyfoot[RE,RO]{
  	\textcolor{white}{%
    	\llap{%
      	\colorbox{blueZUT}{%
        	\makebox[7ex][l]{\sffamily\thepage}%
      }%
      \hspace{-6\marginparsep}%separation from the margin
      \hspace{-5\fboxsep}%
     }%
     }%
    } 
   \addtolength{\headheight}{18pt} % Increase the spacing around the header slightly
}


\renewcommand*{\headrulewidth}{0pt}
\renewcommand*{\footrulewidth}{0pt}

% Removes the header from odd empty pages at the end of chapters
\makeatletter
\renewcommand{\cleardoublepage}{
\clearpage\ifodd\c@page\else
\hbox{}
\vspace*{\fill}
\thispagestyle{empty}
\newpage
\fi}

%----------------------------------------------------------------------------------------
%	BIBLIOGRAPHY (JK - configuring and styling)
%----------------------------------------------------------------------------------------
% 
\usepackage[
style=numeric,% style alphabetic or numeric
citestyle=numeric,
sorting=nyt,%name -year -title
sortcites=true,
autopunct=true,
autolang=hyphen,
hyperref=true, % if the citation is the link to bibliography
backend=biber,
defernumbers=true]{biblatex}
\addbibresource{bibliography.bib} % BibTeX bibliography file
\defbibheading{bibempty}{}


%\usepackage{calc} % For simpler calculation - used for spacing the index letter headings correctly

%----------------------------------------------------------------------------------------
%	CHAPTER & SECTION HEADINGS
%----------------------------------------------------------------------------------------

\usepackage[explicit]{titlesec}

% \titleformat{<command>}[<shape>]{<format>}{<label>}{<sep>}{<before-code>}[<after-code>]

\titleformat{\chapter}[block]
  {\huge\sffamily\color{blueZUT}}
  {\hspace{- 3ex}{\thechapter.} \hspace{0.5em}{#1\strut}}
  {0pt}
  {}

\titleformat{name = \chapter, numberless}[block]
  {\huge\sffamily\color{blueZUT}}
  {{#1\strut}}
  {0pt}
  {}

%----------------------------------------------------------------------------------------
%	Hanging SECTION NUMBERING (MARGIN)
%----------------------------------------------------------------------------------------

\makeatletter
\renewcommand{\@seccntformat}[1]{\llap{\textcolor{blueZUT}{\csname the#1\endcsname}\hspace{1em}}}                    

\renewcommand{\section}
{\@startsection{section}{1}{\z@}
{-4ex \@plus -1ex \@minus -.4ex}
{1ex \@plus.2ex }
{\normalfont\large\sffamily\bfseries}}

\renewcommand{\subsection}{\@startsection {subsection}{2}{\z@}
{-3ex \@plus -0.1ex \@minus -.4ex}
{0.5ex \@plus.2ex }
{\normalfont\sffamily\bfseries}}

\renewcommand{\subsubsection}{\@startsection {subsubsection}{3}{\z@}
{-2ex \@plus -0.1ex \@minus -.2ex}
{.2ex \@plus.2ex }
{\normalfont\small\sffamily\bfseries}}                        

\renewcommand\paragraph{\@startsection{paragraph}{4}{\z@}
{-2ex \@plus-.2ex \@minus .2ex}
{.1ex}
{\normalfont\small\sffamily\bfseries}}



%----------------------------------------------------------------------------------------
%	MAIN TABLE OF CONTENTS (JK modification: style, indentation, colors,)
%----------------------------------------------------------------------------------------

\usepackage{titletoc} % Required for manipulating the table of contents
\contentsmargin{0cm} % Removes the default margin

% Chapter text styling
\titlecontents{chapter}[1.25cm] % Indentation
{\addvspace{12pt}\large\sffamily} % Spacing and font options for chapters
{\color{blueZUT!60}\contentslabel[\Large\thecontentslabel]{1.25cm}\color{blueZUT}} % JK Chapter number
{\color{blueZUT}}  
{\color{blueZUT!60}\normalsize\;\titlerule*[.5pc]{.}\;\thecontentspage} % Page number

% Section text styling
\titlecontents{section}
	[1.25cm] % Left indentation
	{\addvspace{3pt}\sffamily} % Spacing and font options for sections
	{\contentslabel[\thecontentslabel]{1.25cm}} % Formatting of numbered sections of this type
	{} % Formatting of numberless sections of this type
	{\;\titlerule*[.5pc]{.}\;\color{black}\thecontentspage} % Formatting of the filler to the right of the heading and the page number

% Subsection text styling
\titlecontents{subsection}
	[2.5cm] % Left indentation
	{\addvspace{1pt}\sffamily\small} % Spacing and font options for subsections
	{\contentslabel[\thecontentslabel]{1.25cm}} % Formatting of numbered sections of this type
	{} % Formatting of numberless sections of this type
	{\ \titlerule*[.5pc]{.}\;\thecontentspage} % Formatting of the filler to the right of the heading and the page number

%%%%% JK Add figures and tables
% Figure text styling
\titlecontents{figure}
	[0em] % Left indentation
	{\addvspace{3pt}\sffamily} % Spacing and font options for figures
	{\contentslabel[\thecontentslabel]{1.25cm}} % Formatting of numbered sections of this type
	{} % Formatting of numberless sections of this type
	{\ \titlerule*[.5pc]{.}\;\thecontentspage} % Formatting of the filler to the right of the heading and the page number

% Table text styling
\titlecontents{table}
	[0em] % Left indentation
	{\addvspace{3pt}\sffamily} % Spacing and font options for tables
	{\contentslabel[\thecontentslabel]{1.25cm}} % Formatting of numbered sections of this type
	{} % Formatting of numberless sections of this type
	{\ \titlerule*[.5pc]{.}\;\thecontentspage} % Formatting of the filler to the right of the heading and the page number


%----------------------------------------------------------------------------------------
%	THEOREM STYLES
%----------------------------------------------------------------------------------------

\newcommand{\intoo}[2]{\mathopen{]}#1\,;#2\mathclose{[}}
\newcommand{\ud}{\mathop{\mathrm{{}d}}\mathopen{}}
\newcommand{\intff}[2]{\mathopen{[}#1\,;#2\mathclose{]}}
\renewcommand{\qedsymbol}{$\blacksquare$}
\renewcommand{\thmname}{Twierdzenie}

% Boxed/framed environments
\newtheoremstyle{blueZUTnumbox}% Theorem style name
{0pt}% Space above
{0pt}% Space below
{\normalfont}% Body font
{}% Indent amount
{\small\bf\sffamily\color{blueZUT}}% Theorem head font
{\;}% Punctuation after theorem head
{0.25em}% Space after theorem head
{\small\sffamily\color{blueZUT}\thmname{#1}\nobreakspace\thmnumber{\@ifnotempty{#1}{}\@upn{#2}}% Theorem text (e.g. Theorem 2.1)
\thmnote{\nobreakspace\the\thm@notefont\sffamily\bfseries\color{black}---\nobreakspace#3.}} % Optional theorem note

% \newtheoremstyle{blacknumex}% Theorem style name
% {5pt}% Space above
% {5pt}% Space below
% {\normalfont}% Body font
% {} % Indent amount
% {\small\bf\sffamily}% Theorem head font
% {\;}% Punctuation after theorem head
% {0.25em}% Space after theorem head
% {\small\sffamily{\tiny\ensuremath{\blacksquare}}\nobreakspace\thmname{#1}\nobreakspace\thmnumber{\@ifnotempty{#1}{}\@upn{#2}}% Theorem text (e.g. Theorem 2.1)
% \thmnote{\nobreakspace\the\thm@notefont\sffamily\bfseries---\nobreakspace#3.}}% Optional theorem note

\newtheoremstyle{blacknumbox} % Theorem style name
{5pt}% Space above
{5pt}% Space below
{\normalfont}% Body font
{}% Indent amount
{\small\bf\sffamily}% Theorem head font
{\;}% Punctuation after theorem head
{0.25em}% Space after theorem head
{\small\sffamily\thmname{#1}\nobreakspace\thmnumber{\@ifnotempty{#1}{}\@upn{#2}}% Theorem text (e.g. Theorem 2.1)
\thmnote{\nobreakspace\the\thm@notefont\sffamily\bfseries---\nobreakspace#3.}}% Optional theorem note

% Non-boxed/non-framed environments
\newtheoremstyle{blueZUTnum}% Theorem style name
{5pt}% Space above
{5pt}% Space below
{\normalfont}% Body font
{}% Indent amount
{\small\bf\sffamily\color{blueZUT}}% Theorem head font
{\;}% Punctuation after theorem head
{0.25em}% Space after theorem head
{\small\sffamily\color{blueZUT}\thmname{#1}\nobreakspace\thmnumber{\@ifnotempty{#1}{}\@upn{#2}}% Theorem text (e.g. Theorem 2.1)
\thmnote{\nobreakspace\the\thm@notefont\sffamily\bfseries\color{black}---\nobreakspace#3.}} % Optional theorem note
\makeatother

% Defines the theorem text style for each type of theorem to one of the three styles above
\newcounter{dummy} 
\numberwithin{dummy}{section}
\theoremstyle{blueZUTnumbox}
\newtheorem{theoremeT}[dummy]{Twierdzenie}
\theoremstyle{blueZUTnum}
\newtheorem{exampleT}{Przykład}[chapter]
\theoremstyle{blacknumbox}
\newtheorem{definitionT}{Definicja}[section]


%----------------------------------------------------------------------------------------
%	DEFINITION OF COLORED BOXES
%----------------------------------------------------------------------------------------

\RequirePackage[framemethod=default]{mdframed} % Required for creating the theorem, definition, exercise and corollary boxes

% Theorem box
\newmdenv[skipabove=7pt,
skipbelow=7pt,
backgroundcolor=black!3,
linecolor=blueZUT,
innerleftmargin=5pt,
innerrightmargin=5pt,
innertopmargin=5pt,
leftmargin=0cm,
rightmargin=0cm,
innerbottommargin=5pt]{tBox}


% Definition box
\newmdenv[skipabove=7pt,
skipbelow=7pt,
rightline=false,
leftline=true,
topline=false,
bottomline=false,
linecolor=blueZUT,
innerleftmargin=5pt,
innerrightmargin=5pt,
innertopmargin=0pt,
leftmargin=0cm,
rightmargin=0cm,
linewidth=2pt,
innerbottommargin=0pt]{dBox}	

% Creates an environment for each type of theorem and assigns it a theorem text style from the "Theorem Styles" section above and a colored box from above
\newenvironment{theorem}{\begin{tBox}\begin{theoremeT}}{\end{theoremeT}\end{tBox}}				  
\newenvironment{definition}{\begin{dBox}\begin{definitionT}}{\end{definitionT}\end{dBox}}	
\newenvironment{example}{\begin{exampleT}}{\hfill{\tiny\ensuremath{\blacksquare}}\end{exampleT}}		



%----------------------------------------------------------------------------------------
%	LISTING ENVIRONMENT
%----------------------------------------------------------------------------------------

\usepackage{listings}

%Polish set of letteres accepted in the listings
\lstset{
literate=%
{ą}{{\k{a}}}1
{Ą}{{\k{A}}}1
{ć}{{\'c}}1
{Ć}{{\'{C}}}1
{ę}{{\k{e}}}1
{Ę}{{\k{E}}}1
{ł}{{\l{}}}1
{Ł}{{\L{}}}1
{ń}{{\'n}}1
{Ń}{{\'N}}1
{ó}{{\'o}}1
{Ó}{{\'O}}1
{ś}{{\'s}}1
{Ś}{{\'S}}1
{ż}{{\.z}}1
{Ż}{{\.Z}}1
{ź}{{\'z}}1
{Ź}{{\'Z}}1
}

\renewcommand{\lstlistingname}{\small\sffamily\bfseries\color{blueZUT} Algorytm} % Change default listing caption to Algorthm
\renewcommand{\lstlistlistingname}{Lista \lstlistingname ów}

\definecolor{codegreen}{rgb}{0,0.6,0}
\definecolor{codegray}{rgb}{0.5,0.5,0.5}

\lstdefinestyle{mystyle}{
 %   backgroundcolor=\color{grayZUT!10},   
    basicstyle= \small\fontfamily{lmss}\selectfont,%\footnotesize\fontfamily{cmss}\selectfont,
    commentstyle=\color{codegray},
    keywordstyle=\color{violet},
    numberstyle=\tiny\color{codegray},%numeracja linijek
    identifierstyle={\color{black}},
    numbers=left,%numeracja linijek
    numbersep=10pt,%numeracja linijek
    stringstyle=\color{codegreen},
    breakatwhitespace=true,         
    breaklines=true,                 
    captionpos=b,                    
    %keepspaces=false,                
    %showspaces=false,                
    showstringspaces=false,
    showtabs=true,                  
    tabsize=2,
    frame=leftline,
    rulecolor = \color{blueWI},
    xleftmargin=5ex,
    xrightmargin=5ex
    }
 
\lstset{style=mystyle}

%----------------------------------------------------------------------------------------
% CAPTIONS ( JK - design and implementation)
%----------------------------------------------------------------------------------------

\usepackage{caption}
\captionsetup[figure]{name={\small\sffamily\color{blueZUT} Rysunek}}
\captionsetup[table]{name={\small\sffamily\color{blueZUT} Tabela}}
\captionsetup{font={small,sf,singlespacing}}


%----------------------------------------------------------------------------------------
%	HYPERLINKS IN THE DOCUMENTS
%----------------------------------------------------------------------------------------

\usepackage{hyperref}
%\hypersetup{hidelinks,backref=true,pagebackref=true,hyperindex=true,colorlinks=false,breaklinks=true,urlcolor=blueZUT,bookmarks=true,bookmarksopen=false}
\hypersetup{hidelinks,breaklinks=true,urlcolor=blueZUT,bookmarksopen=false,pdftitle={Title},pdfauthor={Author}}


 % JK - Plik zawierający podstawowe elementy konfigurujące układ dokumentu
% UWAGA! - raczej nie będzie potrzeby zmieniania jego struktury

%%%%%%%%%%%%%%%%%%%%%%%%%%%%%%%%%%%%%%%%%
% Plik z definicjami
% Szablon pracy dyplomowej
% Wydział Informatyki 
% Zachodniopomorski Uniwersytet Technologiczny w Szczecinie
% autor Joanna Kołodziejczyk (jkolodziejczyk@zut.edu.pl)
% Bardzo wczesnym pierwowzorem szablonu był
% The Legrand Orange Book
% Version 2.1 (26/09/2018)
%
% Modifications to LOB assigned by %JK
%%%%%%%%%%%%%%%%%%%%%%%%%%%%%%%%%%%%%%%%%



\def\HRule{\color{blueWI} \rule{\linewidth}{0.6pt}} % horisontal rule in ZUT color

%----------------------------------------------------------------------------------------
% Typ pracy (wybrać właściwy)
%----------------------------------------------------------------------------------------
%\def\degreename{Praca dyplomowa magisterska} 
\def\degreename{Praca dyplomowa inżynierska}
%\def\degreename{Praca dyplomowa licencjacka}

%----------------------------------------------------------------------------------------
% Temat pracy
%----------------------------------------------------------------------------------------
\def\ttitle{	Aplikacja porównująca działanie klasyfikatorów zespołowych z innymi wybranymi klasyfikatorami}
\def\ttitleEng{	Application comparing ensemble classifiers with other selected classifiers} %temat pracy w j. angielskim


%----------------------------------------------------------------------------------------
% Informacje o autorze
%----------------------------------------------------------------------------------------
\def\authornames{Marcin Mielniczuk} %imię i nazwisko autora
\def\noalbum{51621}
\def\speciality{Inżynieria oprogramowania} %nazwa specjalności
\def\field{Informatyka} %dziedzina nauki

%----------------------------------------------------------------------------------------
% Informacje o promotorze
%----------------------------------------------------------------------------------------
\def\supname{dr inż. Aneta Bera} %imię i nazwisko promotora
\def\departmentname{Katedra Metod Sztucznej Inteligencji i Matematyki Stosowanej} %nazwa katedry promotora

%----------------------------------------------------------------------------------------
% Informacje o promotorze dodatkowym (odkomentować w razie potrzeby)
%----------------------------------------------------------------------------------------
% \def\supExt{prof. dr hab. inż. Marian Kowalski} %imię i nazwisko promotora
% \def\departmentnameExt{Katedra Grafiki Komputerowej} %nazwa katedry promotora

%----------------------------------------------------------------------------------------
% Data wydania tematu pracy
%----------------------------------------------------------------------------------------
\def\datetitle{20.01.2025}

%----------------------------------------------------------------------------------------
% Rok i miejsce złożenia pracy
%----------------------------------------------------------------------------------------
\def\placesubmit{Szczecin}
\def\yearsubmit{202X}
  % JK - dodatkowe definicje głównie treść strony tytułowej
% UWAGA! - konieczność edycji celem zmiany autora/tematu/dat itp., itd

%----------------------------------------------------------------------------------------
% OTWARCIE DOKUMENTU
%----------------------------------------------------------------------------------------
\begin{document}

%----------------------------------------------------------------------------------------
% STRONA TYTUŁOWA 
%----------------------------------------------------------------------------------------
%%%%%%%%%%%%%%%%%%%%%%%%%%%%%%%%%%%%%%%%%
% Układ strony tytułowej
% Szablon pracy dyplomowej
% Wydział Informatyki 
% Zachodniopomorski Uniwersytet Technologiczny w Szczecinie
% autor Joanna Kołodziejczyk (jkolodziejczyk@zut.edu.pl)
% Bardzo wczesnym pierwowzorem szablonu był
% The Legrand Orange Book
% Version 2.1 (26/09/2018)
%
% Modifications to LOB assigned by %JK
%%%%%%%%%%%%%%%%%%%%%%%%%%%%%%%%%%%%%%%%%

%----------------------------------------------------------------------------------------
%	STRONA TYTUŁOWA - NIE ZMIENIAĆ
%----------------------------------------------------------------------------------------
\begingroup
\firaoldstyle  %Set sans serif fonts for title page
\centering %Center all paragraphs
\thispagestyle{empty} % Suppress headers and footers on the title page

%----------------------------------------------------------------------------------------
%	LOGOTYPY - NIE ZMIENIAĆ
%----------------------------------------------------------------------------------------
    \includegraphics[scale=.4]{logo5.jpg}\\[.75cm]

%%%
Wydział Informatyki\\[.25cm]
kierunek studiów: \field\\[.25cm]
specjalność: \speciality\\[1.75cm]
%----------------------------------------------------------------------------------------
%	TYP - NIE ZMIENIAĆ
%----------------------------------------------------------------------------------------
{\Large \degreename}\\[1cm]

%----------------------------------------------------------------------------------------
%	TYTUŁ - NIE ZMIENIAĆ
%----------------------------------------------------------------------------------------
% {\color{blueZUT} 
\Large {\bfseries \MakeUppercase  \ttitle }\\[.5cm]% Thesis title
\large {\bfseries \MakeUppercase \ttitleEng} \\[.5cm]% Thesis title in English

%----------------------------------------------------------------------------------------
%	AUTOR - NIE ZMIENIAĆ
%----------------------------------------------------------------------------------------
 % {\color{blueZUT} {
{\large \bfseries \authornames}\\[.5cm]
{nr albumu: \bfseries {\noalbum}}\\[1cm]

%----------------------------------------------------------------------------------------
%	OPIEKUN PRACY - NIE ZMIENIAĆ
%----------------------------------------------------------------------------------------

{\large  Opiekun:}\\[.5cm]
{\large \bfseries{\supname}}\\[.5cm]
{\departmentname}\\[1cm]

%----------------------------------------------------------------------------------------
%	DODATKOWY OPIEKUN PRACY - NIE ZMIENIAĆ
%----------------------------------------------------------------------------------------

\ifdefined\supExt
    {\large Opiekun zewnętrzny:}\\[.5cm]
    {\large \bfseries{\supExt}}\\[.5cm]
    \departmentnameExt
\fi

%----------------------------------------------------------------------------------------
%	DÓŁ STRONY - NIE ZMIENIAĆ
%----------------------------------------------------------------------------------------
~\vfill
\placesubmit, \yearsubmit\\

\endgroup



%----------------------------------------------------------------------------------------
% PUSTA STRONA PO STRONIE TYUŁOWEJ (JK - design and implementation)
%----------------------------------------------------------------------------------------
\newpage
\thispagestyle{empty}
~\vfill

%----------------------------------------------------------------------------------------
%	STRESZCZENIE I SŁOWA KLUCZOWE (1 STRONA) (JK - design and implementation)
%----------------------------------------------------------------------------------------
\newpage
\thispagestyle{empty}
%%%%%%%%%%%%%%%%%%%%%%%%%%%%%%%%%%%%%%%%%
% Specjalna strona pracy ze streszczeniem i abstractem w j. angielskim
% Szablon pracy dyplomowej
% Wydział Informatyki 
% Zachodniopomorski Uniwersytet Technologiczny w Szczecinie
% autor Joanna Kołodziejczyk (jkolodziejczyk@zut.edu.pl)
% Bardzo wczesnym pierwowzorem szablonu był
% The Legrand Orange Book
% Version 2.1 (26/09/2018)
%
% Modifications to LOB assigned by %JK
%%%%%%%%%%%%%%%%%%%%%%%%%%%%%%%%%%%%%%%%%


\begin{center}
\noindent {{\color{blueZUT}\Large\sffamily  {Streszczenie}}}\\[1cm] 
\end{center}

W tym miejscu trzeba napisać streszczenie pracy w języku polskim. Zawiera krótką charakterystykę dziedziny, przedmiotu i wyników zaprezentowanych w pracy. Maksymalnie 1/2 strony.

\vspace{10pt}
\noindent{\bf słowa kluczowe:} np. informatyka, sterowanie, grafika komputerowa

\vfill

\begin{center}
\noindent {{\color{blueZUT}\Large\sffamily {Abstract}}}\\[1cm] 
\end{center}
The abstract's purpose, which should not exceed 150 words, is to provide sufficient information to allow potential readers to decide on the thesis's relevance—a maximum of half the page.

\vspace{10pt}
\noindent{\bf keywords:} e.g.: computer science, control, computer graphics %abstract.tex 
 
%----------------------------------------------------------------------------------------
%	SPIS TREŚCI (JK - design and implementation)
%----------------------------------------------------------------------------------------
\pagestyle{empty}  % Wyłącz stopkę i nagłówek w TOC
\tableofcontents % wyświetl spis
\addtocontents{toc}{\protect\thispagestyle{empty}} % Zachowaj pusty nagłówek i stopkę w TOC
% Wymuszenie rozpoczęcie pierwszego rozdziału na nieparzystej stronie, aby znajdował się po prawej stronie 
\cleardoublepage 
% Ponownie włącz nagłówki i stopki
\pagestyle{fancy} 

%----------------------------------------------------------------------------------------
%	WSTĘP W PRACY DYPLOMOWEJ (JK - design and implementation)
%----------------------------------------------------------------------------------------
\addcontentsline{toc}{chapter}{Wstęp}
%%%%%%%%%%%%%%%%%%%%%%%%%%%%%%%%%%%%%%%%%
% Plik z wstępem do pracy
% Szablon pracy dyplomowej
% Wydział Informatyki 
% Zachodniopomorski Uniwersytet Technologiczny w Szczecinie
% autor Joanna Kołodziejczyk (jkolodziejczyk@zut.edu.pl)
% Bardzo wczesnym pierwowzorem szablonu był
% The Legrand Orange Book
% Version 2.1 (26/09/2018)
%
% Modifications to LOB assigned by %JK
%%%%%%%%%%%%%%%%%%%%%%%%%%%%%%%%%%%%%%%%%


\chapter*{Wstęp}

Wstęp powinien być nie dłuższy niż 2 strony. Najlepiej napisać go dopiero, gdy praca jest już skończona i wszystkie jej części spisane.

Wstęp powinien zawierać:

\begin{enumerate}
\item Opis dziedziny jakiej dotyczy praca, ze wskazaniem, że temat pracy jest ważny, bieżący, itp.
\item Jaki problem z dziedziny się rozwiązuje.
\item Cel i teza pracy
\item W jaki sposób cel zostanie osiągnięty a tez potwierdzona.
\item Struktura pracy.
\end{enumerate}  % introduction.tex zawiera treść wstępu

%----------------------------------------------------------------------------------------
%	ROZDZIAŁ 1
%----------------------------------------------------------------------------------------
%%%%%%%%%%%%%%%%%%%%%%%%%%%%%%%%%%%%%%%%%
% Szablon pracy dyplomowej
% Wydział Informatyki 
% Zachodniopomorski Uniwersytet Technologiczny w Szczecinie
% autor Joanna Kołodziejczyk (jkolodziejczyk@zut.edu.pl)
% Bardzo wczesnym pierwowzorem szablonu był
% The Legrand Orange Book
% Version 2.1 (26/09/2018)
%
% Modifications to LOB assigned by %JK
%%%%%%%%%%%%%%%%%%%%%%%%%%%%%%%%%%%%%%%%%

%----------------------------------------------------------------------------------------
%	CHAPTER 1
% 	author: Joanna Kolodziejczyk (jkolodziejczyk@zut.edu.pl)
%----------------------------------------------------------------------------------------

\chapter{Informacje podstawowe dotyczące szablonu}
\label{rozdzial1}

Niniejszy szablon przeznaczony jest do utrzymania spójnego wyglądu prac dyplomowych na Wydziale Informatyki Zachodniopomorskiego Uniwersytetu Technologicznego w Szczecinie. Na bazie niniejszego szablonu przygotowywane są podręczniki, skrypty i monografie wydziałowe.

\section{Autorzy stylu}
\index{Autorzy stylu}

Niniejszy szablon powstał na bazie szablonu\footnote{\url{http://www.LaTeXTemplates.com}} o nazwie ,,The Legrand Orange Book''. 
Autorzy pierwotnego stylu to:
\begin{itemize}
\item Mathias Legrand (legrand.mathias@gmail.com)
\item Vel (vel@latextemplates.com).
\end{itemize}

Późniejsze zmiany na potrzeby stylu Wydziału Informatyki zostały wprowadzone przez Joannę Kołodziejczyk (jkolodziejczyk@zut.edu.pl). Wizualna strona stylu  jest konsensusem uzyskanym w wyniku prac zespołu: 
\begin{itemize}
\item Anna Barcz
\item Marcin Korzeń
\item Mirosław Łazoryszczak
\item Piotr Piela
\end{itemize}

Niniejszy styl jest udostępniony na zasadach licencji CC BY-NC-SA 3.0 ({\url{http://creativecommons.org/licenses/by-nc-sa/3.0/}}).

\section{Kompilowanie dokumentu}
\index{Kompilacja}

 Do kompilacji wykorzystać należy procesor {\em pdflatex}.
 
 Ten szablon wykorzystuje pakiet - menedżer spisu literatury o nazwie {\em biber}\footnote{ \url{http://biblatex-biber.sourceforge.net}}, który jest dostępny w wielu popularnych środowiskach programistycznych dla \LaTeX.  

Najpewniejszym sposobem na poprawną kompilację jest wykonanie jej z wiersza poleceń następującą sekwencją rozkazów:

\begin{lstlisting}[language=bash, caption=Skrypt kompilujący, label=alg:1]
pdflatex main
biber main
pdflatex main
pdflatex main

\end{lstlisting}

Szablon wykorzystuje również rozliczne pakiety, które mogą wymagać
zaktualizowane do najnowszych wersji. Rekomendowane jest więc, by utrzymywać swoją wersję dystrybucji \LaTeX w najnowszej wersji.

\section{Zawartość pakietu}\index{Zawartość pakietu}

Pakiet ze stylem zawiera pliki tekstowe z zawartością przykładowej pracy , czyli w tym wypadku niniejszego przewodnika oraz pliki konfiguracyjne. W tabeli \ref{tab:1} znajduje się lista plików wraz z ich krótkim opisem.

\begin{table}[h]
\centering
\caption{Pliki w pakiecie szablonu}
\begin{tabular}{p{.2\textwidth}|p{.1\textwidth} | p{.60\textwidth}}
\toprule
\textbf{Nazwa pliku}& \textbf{Edycja}  & \textbf{Opis}\\
\midrule
main.tex &TAK& Główny plik pracy. Zawiera otwarcie dokumentu, dołącza pliki konfiguracyjne, spis treści, dołącza wszystkie rozdziały, literaturę oraz zamyka dokument.  \\
definitions.tex &TAK& Zawiera elementy strony tytułowej: autora, tytuł, niezbędne informacje ze strony tytułowej.\\
abstract.tex &TAK& Zawartość streszczenia pracy w j. polskim i j. angielskim. Zaleca się, by nie przekroczyć tekstem 1 strony.\\
introduction.tex&TAK & Wstęp pracy dyplomowej\\
chapter1.tex &TAK& Zawartość rozdziału pierwszego.  \\
chapter2.tex  &TAK& Zawartość rozdziału drugiego.  \\
conclusions.tex &TAK& Zakończenie/podsumowanie pracy dyplomowej.\\
bibliography.bib &TAK & Plik z literaturą w formacie bib.  \\
Pictures  &TAK & Katalog zawierający ilustracje i  grafiki z pracy. Nazwa katalogu jest wykorzystana w pliku konfiguracyjnym jako katalog domyślny do każdego osadzonego obiektu graficznego.  \\
title\_page.tex&NIE & Struktura strony tytułowej.\\
structure.tex  &NIE& Plik konfiguracyjny (preambuła). Podłącza i konfiguruje wszystkie niezbędne pakiety oraz definiuje otoczenia wykorzystywane w szablonie. \\
\bottomrule
\end{tabular}
\label{tab:1} 
\end{table}
%------------------------------------------------

\subsection{Podstawowe wymagania dotyczące szablonu}

Praca ma ustawiony format wielkości strony A4, w pliku {\em structure.tex} w konfiguracji pakietu {\bf geometry}. Przygotowana jest do druku dwustronnego. Ustawiony jest 1cm margines na oprawę.

\noindent{\bf UWAGA!} Każdy rozdział zaczyna się na stronie nieparzystej


\subsection{Bibliografia}

Po zaktualizowaniu pliku z literaturą należy dokonać kompilacji z użyciem procesora {\em biber} a nie standardowego  {\em Bibtex}. Następnie dwukrotnie skompilować z użyciem {\em pdflatex}, aby zmiany były widoczne. Pakiet do zarządzania bibliografią to {\em biblatex}. Jego konfiguracja jest dostępna w pliku {\em structure.tex}

Bibliografia w pracach dyplomowych ma format numeryczny, sortowanie alfabetyczne i jest podzielona na  części: książki, artykuły i źródła internetowe i inne. Konfiguracja jest zapisana w niniejszym szablonie.

Tylko zacytowane prace znajdą się w spisie literatury. Nie należy tego zmieniać. Cytowanie odbywa się przez użycie polecenia $\backslash cite\{\}$ a efekt jest następujący: \cite{book_key}, \cite{article_key}, \cite{lam1994}, \cite{knuthwebsite}.





 % chapter1.tex zawiera treść rozdziału 1

%----------------------------------------------------------------------------------------
%	ROZDZIAŁ 2
%----------------------------------------------------------------------------------------
%%%%%%%%%%%%%%%%%%%%%%%%%%%%%%%%%%%%%%%%%
% Szablon pracy dyplomowej
% Wydział Informatyki 
% Zachodniopomorski Uniwersytet Technologiczny w Szczecinie
% autor Joanna Kołodziejczyk (jkolodziejczyk@zut.edu.pl)
% Bardzo wczesnym pierwowzorem szablonu był
% The Legrand Orange Book
% Version 2.1 (26/09/2018)
%
% Modifications to LOB assigned by %JK
%%%%%%%%%%%%%%%%%%%%%%%%%%%%%%%%%%%%%%%%%


%----------------------------------------------------------------------------------------
%	CHAPTER 2
%----------------------------------------------------------------------------------------

\chapter{Przedstawienie zestawu zdefiniowanych otoczeń możliwych do wykorzystania w pracy}
\chaptermark{Elementy w pracy} % Tekst, który wyświetli się w nagłówku strony,  jeżeli jest za długi tytuł rozdziału

Szablon definiuje różne otoczenia, z których może korzystać autor pracy dyplomowej. Takowe zostąły opisane w niniejszym rozdziale. Jeżeli otoczenie jest zawarte w szablonie, to właśnie takie należy stosować. Jeżeli w ramach pracy zachodzi potrzeba korzystania z elementów nie ujętych w szablonie można je stosować dowolnie.

\section{Listy}

Strukturę list wyliczeniowych poprawnie wykorzystywanych w języku polskim opisują słowniki. Oto fragment z poradnika \url{sjp.pwn.pl}. Podpunkty zdaniowe powinny się zaczynać wielką literą, a niezdaniowe małą. Jeśli podpunkt jest zdaniem, to zamykamy go kropką lub znakiem równoważnym, jeśli nie – przecinkiem lub średnikiem.  Nagłówek listy powinien się kończyć dwukropkiem nawet wtedy, gdy następują po nim podpunkty zdaniowe, pisane wielką literą. 

\noindent Listy numerowane wykorzystuje się, gdy kolejność elementów jest istotna:
\begin{enumerate}
\item element 1,
\item element 2,
\item element 34234.
\end{enumerate}

\noindent Listy wypunktowane wykorzystuje się, gdy kolejność elementów jest istotna:
\begin{itemize}
\item osoba Alicja A, 
\item zwierzak kot Alicji,
\item dom Alicji.
\end{itemize}

Jeżeli zamiast punktora lub wyliczenia istnieje potrzeba wykorzystania innego słowa jako wyliczenia,to możliwe jest stosowanie listy z nagłówkami (descriptions-definitions),
\begin{description}
\item[Nazwa] Opis;
\item[Słowo] Definicja;
\item[Komentarz] Wywód.
\end{description}

\section{Tabele i rysunki w pracy}

Tabele podpisuje się z góry. Ważne, że każdy element osadzony w treści pracy musi zostać przywołany w tekście np. w Tabeli: \ref{tab:eksperyment1} zaprezentowano wyniki eksperymentu

\begin{table}[h]
\centering
\caption{Tabelka z wynikami eksperymentu 1} 
\begin{tabular}{l l l}
\toprule
\textbf{Lek} & \textbf{Odpowiedź 1} & \textbf{Odpowiedź 2}\\
\midrule
Lek 1 & 0.0003262 & 0.562 \\
Lek 2 & 0.0015681 & 0.910 \\
Lek 3 & 0.0009271 & 0.296 \\
\bottomrule
\end{tabular}
\label{tab:eksperyment1} % Etykiety dla każdej tabeli musi być unikalna
\end{table}

Rysunki podpisywane są z dołu. Należy tak jak w tabeli w treści pracy odnieść się do etykiety rysunku, np. Rysunek \ref{fig:rys1} przedstawia

\begin{figure}[!ht]
		\centering
		%\vspace{2mm}
%		\includegraphics[scale=0.6]{Dyplom-styl/Pictures/placeholder.jpg}
		\caption{To jest przykład osadzania rysunków (źródło:)}
		\label{fig:rys1}
\end{figure}



\section{Twierdzenia i definicje}

W pracach o rozbudowanej części teoretycznej może istnieć potrzeba prezentowania twierdzeń i definicji. Poniżej przykłady stosowania odpowiednich otoczeń. 

A oto twierdzenie składające się z kilku równań.

\begin{theorem}[Tytuł twierdzenia]
In $E=\mathbb{R}^n$ all norms are equivalent. It has the properties:
\begin{align}
& \big| ||\mathbf{x}|| - ||\mathbf{y}|| \big|\leq || \mathbf{x}- \mathbf{y}||\\
&  ||\sum_{i=1}^n\mathbf{x}_i||\leq \sum_{i=1}^n||\mathbf{x}_i||\quad\text{where $n$ is a finite integer}
\end{align}
\end{theorem}

Jest to twierdzenie składające się z tylko jednej linii.

\begin{theorem}
Zbiór $\mathcal{D}(G)$ ma gęstość $L^2(G)$, $|\cdot|_0$. 
\end{theorem}


To jest przykład definicji. Definicja może być matematyczna lub może definiować koncepcję. Jeżeli w pracy wprowadza się nowe pojęcia najlepiej ująć je w formie definicji, co będzie spójnie akcentować taki fakt.

\begin{definition}[Tytuł definicji]
Given a vector space $E$, a norm on $E$ is an application, denoted $||\cdot||$, $E$ in $\mathbb{R}^+=[0,+\infty[$ such that:
\begin{align}
& ||\mathbf{x}||=0\ \Rightarrow\ \mathbf{x}=\mathbf{0}\\
& ||\lambda \mathbf{x}||=|\lambda|\cdot ||\mathbf{x}||\\
& ||\mathbf{x}+\mathbf{y}||\leq ||\mathbf{x}||+||\mathbf{y}||
\end{align}
\end{definition}


%------------------------------------------------

\section{Przykłady}

W wielu pracach do przedstawionych idei prezentuje się przykłady. Na tą okoliczność przygotowano otoczenie {\it example}, które pomoże utrzymać spójny wygląd wszystkich przykładów. 

Przykład może dotyczyć obliczeń lub prezentacji działania np. metody, algorytmu kodu, itp, itd.

\begin{example}

Let $G=\{x\in\mathbb{R}^2:|x|<3\}$ and denoted by: $x^0=(1,1)$; consider the function:
\begin{equation}
f(x)=\left\{\begin{aligned} & \mathrm{e}^{|x|} & & \text{si $|x-x^0|\leq 1/2$}\\
& 0 & & \text{si $|x-x^0|> 1/2$}\end{aligned}\right.
\end{equation}
The function $f$ has bounded support, we can take $A=\{x\in\mathbb{R}^2:|x-x^0|\leq 1/2+\epsilon\}$ for all $\epsilon\in\intoo{0}{5/2-\sqrt{2}}$.
\end{example}


\begin{example}[Przykład krzyżowania]
\begin{align*}
&\text{Pokolenie rodziców:} & & & &\text{Pokolenie potomków:} \\
&x_1 = (\textbf{00110} | \textbf{011}) & &\xrightarrow{\text{krzyżowanie}} & &(\textbf{00110}|1 0 1)\\
&x_2 = (0 1 1 0 1 | 1 0 1) & & & &(0 1 1 0 1|\textbf{011})\\
\end{align*}
\end{example}


%------------------------------------------------

\section{Kod źródłowy}

Jeżeli w pracy prezentowany jest kod źródłowy, to należy skorzystać z predefiniowanego i skonfigurowanego (w structure.tex) otoczenia lstlisting. Jeżeli autor chciałby umieścić w pracy pseudokod, to może posłużyć się wybranym przez siebie otoczeniem.

Zasada odwołania w tekście do algorytmu jest taka sama jak do tabeli i rysunku np. algorytm \ref{alg:kod1} został opracowany na potrzeby...

\begin{lstlisting}[language=Python, caption=Fragment algorytmu xxx, label=alg:kod1]
import numpy as np
# change this value for a different result
celsius = 37.5

# calculate fahrenheit
fahrenheit = (celsius * 1.8) + 32
print('%0.1f degree Celsius is equal to %0.1f degree Fahrenheit' %(celsius,fahrenheit))
 
def incmatrix(genl1,genl2):
    m = len(genl1)
    n = len(genl2)
    M = None #to become the incidence matrix
    VT = np.zeros((n*m,1), int)  #zmienna ąźćł
 
    #compute the bitwise xor matrix
    M1 = bitxormatrix(genl1)
    M2 = np.triu(bitxormatrix(genl2),1) 
 
    for i in range(m-1):
        for j in range(i+1, m):
            [r,c] = np.where(M2 == M1[i,j])
            for k in range(len(r)):
                VT[(i)*n + r[k]] = 1;
                VT[(i)*n + c[k]] = 1;
                VT[(j)*n + r[k]] = 1;
                VT[(j)*n + c[k]] = 1;
 
                if M is None:
                    M = np.copy(VT)
                else:
                    M = np.concatenate((M, VT), 1)
 
    return M
\end{lstlisting}
 % chapter2.tex zawiera treść rozdziału 2

%----------------------------------------------------------------------------------------
%	ROZDZIAŁy kolejne należy dodać analogicznie do 1 i 2 
% utworzyć pliki i je załączyć (include)
%----------------------------------------------------------------------------------------

%----------------------------------------------------------------------------------------
%	ZAKOŃCZENIE PRACY DYPLOMOWEJ (JK - design and implementation)
%----------------------------------------------------------------------------------------
\addcontentsline{toc}{chapter}{Podsumowanie}
%%%%%%%%%%%%%%%%%%%%%%%%%%%%%%%%%%%%%%%%%
% Wnioski do pracy dyplomowej
% Szablon pracy dyplomowej
% Wydział Informatyki 
% Zachodniopomorski Uniwersytet Technologiczny w Szczecinie
% autor Joanna Kołodziejczyk (jkolodziejczyk@zut.edu.pl)
% Bardzo wczesnym pierwowzorem szablonu był
% The Legrand Orange Book
% Version 2.1 (26/09/2018)
%
% Modifications to LOB assigned by %JK
%%%%%%%%%%%%%%%%%%%%%%%%%%%%%%%%%%%%%%%%%


\chapter*{Podsumowanie}

Podsumowanie pracy powinno na maksymalnie dwóch stronach przedstawić główne wyniki pracy dyplomowej. Struktura zakończenia to:
\begin{enumerate}
\item Przypomnienie celu i hipotez
\item Co w pracy wykonano by cel osiągnąć (analiza, projekt, oprogramowanie, badania eksperymentalne)
\item Omówienie głównych wyników pracy
\item Jak wyniki wzbogacają dziedzinę
\item Zamknięcie np. poprzez wskazanie dalszych kierunków badań.
\end{enumerate} % conclusions.tex zawiera treść zakończenia/podsumowania pracy dyplomowej


%----------------------------------------------------------------------------------------
%	SPIS LITERATURY (JK - design and implementation)
%----------------------------------------------------------------------------------------
\pagestyle{empty} 
\chapter*{Spis literatury}
\addcontentsline{toc}{chapter}{\textcolor{blueZUT}{Spis literatury}}
% Poniżej zdefiniowane są filtry do dzielenia spisu literatury na kategorie
\defbibfilter{articles}{
  type=article or 
  type=inproceedings
}
\defbibfilter{other}{
  type=url or
  type=online  or
  type=manual or
  type=misc
}

% UWAGA! -aby zmienić zawartość spisu literatury należy wyedytować plik
% bibliography.bib

%------------------------------------------------
% Spis literartury podzielony jest na 3 kategorie
% 1
\section*{Książki}
\addcontentsline{toc}{section}{Książki}
\printbibliography[heading=bibempty,type=book]

%------------------------------------------------
% 2
\section*{Artykuły}
\addcontentsline{toc}{section}{Artykuły}
\printbibliography[heading=bibempty,filter=articles]
%------------------------------------------------
% 3
\section*{Źródła internetowe i inne}
\addcontentsline{toc}{section}{Źródła internetowe i inne}
\printbibliography[heading=bibempty,filter=other]


%----------------------------------------------------------------------------------------
%	APPENDIX (JK - design and implementation)
%----------------------------------------------------------------------------------------
\pagestyle{fancy} 
\begin{appendix}
\appendix
%%%%%%%%%%%%%%%%%%%%%%%%%%%%%%%%%%%%%%%%%
% Szablon pracy dyplomowej
% Wydział Informatyki 
% Zachodniopomorski Uniwersytet Technologiczny w Szczecinie
% autor Joanna Kołodziejczyk (jkolodziejczyk@zut.edu.pl)
% Bardzo wczesnym pierwowzorem szablonu był
% The Legrand Orange Book
% Version 2.1 (26/09/2018)
%
% Modifications to LOB assigned by %JK
%%%%%%%%%%%%%%%%%%%%%%%%%%%%%%%%%%%%%%%%%


\chapter{Dodatek}
\label{chapter:dodatek_A}

\textit{W dodatkach mogą znaleźć się większe fragmenty kodów źródłowych, instrukcje działania programów, specyfikacje opcji, z którymi wywołuje się program, większe dane tabelaryczne, specyfikacje oprogramowania, dłuższe opisy teoretyczne, itp.}

What is Lorem Ipsum?
Lorem Ipsum is simply dummy text of the printing and typesetting industry. Lorem Ipsum has been the industry'es standard dummy text ever since the 1500s, when an unknown printer took a galley of type and scrambled it to make a type specimen book. It has survived not only five centuries, but also the leap into electronic typesetting, remaining essentially unchanged. It was popularised in the 1960s with the release of Letraset sheets containing Lorem Ipsum passages, and more recently with desktop publishing software like Aldus PageMaker including versions of Lorem Ipsum.

Why do we use it?
It is a long established fact that a reader will be distracted by the readable content of a page when looking at its layout. The point of using Lorem Ipsum is that it has a more-or-less normal distribution of letters, as opposed to using 'Content here, content here', making it look like readable English. Many desktop publishing packages and web page editors now use Lorem Ipsum as their default model text, and a search for 'lorem ipsum' will uncover many web sites still in their infancy. Various versions have evolved over the years, sometimes by accident, sometimes on purpose (injected humour and the like).


Where does it come from?
Contrary to popular belief, Lorem Ipsum is not simply random text. It has roots in a piece of classical Latin literature from 45 BC, making it over 2000 years old. Richard McClintock, a Latin professor at Hampden-Sydney College in Virginia, looked up one of the more obscure Latin words, consectetur, from a Lorem Ipsum passage, and going through the cites of the word in classical literature, discovered the undoubtable source. Lorem Ipsum comes from sections 1.10.32 and 1.10.33 of "de Finibus Bonorum et Malorum" (The Extremes of Good and Evil) by Cicero, written in 45 BC. This book is a treatise on the theory of ethics, very popular during the Renaissance. The first line of Lorem Ipsum, "Lorem ipsum dolor sit amet..", comes from a line in section 1.10.32.

The standard chunk of Lorem Ipsum used since the 1500s is reproduced below for those interested. Sections 1.10.32 an

What is Lorem Ipsum?
Lorem Ipsum is simply dummy text of the printing and typesetting industry. Lorem Ipsum has been the industry's standard dummy text ever since the 1500s, when an unknown printer took a galley of type and scrambled it to make a type specimen book. It has survived not only five centuries, but also the leap into electronic typesetting, remaining essentially unchanged. It was popularised in the 1960s with the release of Letraset sheets containing Lorem Ipsum passages, and more recently with desktop publishing software like Aldus PageMaker including versions of Lorem Ipsum.

Why do we use it?
It is a long established fact that a reader will be distracted by the readable content of a page when looking at its layout. The point of using Lorem Ipsum is that it has a more-or-less normal distribution of letters, as opposed to using 'Content here, content here', making it look like readable English. Many desktop publishing packages and web page editors now use Lorem Ipsum as their default model text, and a search for 'lorem ipsum' will uncover many web sites still in their infancy. Various versions have evolved over the years, sometimes by accident, sometimes on purpose (injected humour and the like).


Where does it come from?
Contrary to popular belief, Lorem Ipsum is not simply random text. It has roots in a piece of classical Latin literature from 45 BC, making it over 2000 years old. Richard McClintock, a Latin professor at Hampden-Sydney College in Virginia, looked up one of the more obscure Latin words, consectetur, from a Lorem Ipsum passage, and going through the cites of the word in classical literature, discovered the undoubtable source. Lorem Ipsum comes from sections 1.10.32 and 1.10.33 of "de Finibus Bonorum et Malorum" (The Extremes of Good and Evil) by Cicero, written in 45 BC. This book is a treatise on the theory of ethics, very popular during the Renaissance. The first line of Lorem Ipsum, "Lorem ipsum dolor sit amet..", comes from a line in section 1.10.32.

The standard chunk of Lorem Ipsum used since the 1500s is reproduced below for those interested. Sections 1.10.32 anWhat is Lorem Ipsum?
Lorem Ipsum is simply dummy text of the printing and typesetting industry. Lorem Ipsum has been the industry's standard dummy text ever since the 1500s, when an unknown printer took a galley of type and scrambled it to make a type specimen book. It has survived not only five centuries, but also the leap into electronic typesetting, remaining essentially unchanged. It was popularised in the 1960s with the release of Letraset sheets containing Lorem Ipsum passages, and more recently with desktop publishing software like Aldus PageMaker including versions of Lorem Ipsum.

Why do we use it?
It is a long established fact that a reader will be distracted by the readable content of a page when looking at its layout. The point of using Lorem Ipsum is that it has a more-or-less normal distribution of letters, as opposed to using 'Content here, content here', making it look like readable English. Many desktop publishing packages and web page editors now use Lorem Ipsum as their default model text, and a search for 'lorem ipsum' will uncover many web sites still in their infancy. Various versions have evolved over the years, sometimes by accident, sometimes on purpose (injected humour and the like).


Where does it come from?
Contrary to popular belief, Lorem Ipsum is not simply random text. It has roots in a piece of classical Latin literature from 45 BC, making it over 2000 years old. Richard McClintock, a Latin professor at Hampden-Sydney College in Virginia, looked up one of the more obscure Latin words, consectetur, from a Lorem Ipsum passage, and going through the cites of the word in classical literature, discovered the undoubtable source. Lorem Ipsum comes from sections 1.10.32 and 1.10.33 of "de Finibus Bonorum et Malorum" (The Extremes of Good and Evil) by Cicero, written in 45 BC. This book is a treatise on the theory of ethics, very popular during the Renaissance. The first line of Lorem Ipsum, "Lorem ipsum dolor sit amet..", comes from a line in section 1.10.32.

The standard chunk of Lorem Ipsum used since the 1500s is reproduced below for those interested. Sections 1.10.32 an % conclusions.tex zawiera treść zakończenia/podsumowania pracy dyplomowej
\end{appendix}
%----------------------------------------------------------------------------------------
% Koniec dokumentu
\end{document}
